\documentclass[11pt]{article}
\usepackage[margin=1in]{geometry}
\usepackage[T1]{fontenc}
\usepackage{lmodern}
\usepackage{amsmath,amssymb,amsthm,mathtools}
\usepackage{graphicx}
\usepackage{microtype}
\usepackage[hidelinks]{hyperref}

\newtheorem{theorem}{Theorem}
\newtheorem{lemma}[theorem]{Lemma}
\newcommand{\D}{\mathbb D}
\newcommand{\T}{\mathbb T}
\newcommand{\C}{\mathbb C}
\newcommand{\Hh}{\mathcal H}
\newcommand{\spn}{\operatorname{span}}

\title{A Scalar Sharp Kernel Whose Coordinate Multiplier Is Not Weakly Homogeneous}
\author{Candidate full solution packet\\\small arXiv:1811.05428, Section 3, PDF page 8}
\date{17 August 2026}

\begin{document}
\maketitle

\noindent\textbf{Status.}
Candidate full affirmative answer, subject to human review.  The construction
is elementary and appears not to require the matrix-valued jet machinery used
in the source paper.

\section{Source question}

Ghara \cite[Section 3, p.~8]{Ghara} notes that examples are known for
matrix-valued sharp kernels, but leaves open the scalar case:

\begin{center}
\includegraphics[width=\textwidth]{figures/open_problem_crop.png}
\end{center}

In the terminology of the source, a scalar positive-definite kernel $K$ on
$\D$ is \emph{sharp} when the coordinate multiplier $M_z$ is bounded and
\[
 \ker(M_z^*-\overline w I)=\spn\{K(\,\cdot\,,w)\}
 \qquad(w\in\D).
\]
The question asks whether such a kernel can have $M_z$ not weakly
homogeneous.  Recall that weak homogeneity requires
$\varphi(M_z)$ to be similar to $M_z$ for every disk automorphism $\varphi$.

\section{Idea of the proof}

Start with the Hardy space, but adjoin one holomorphic function having a
single boundary pole:
\[
 g(z)=\frac1{1-z},\qquad \Hh=H^2\mathbin{\dotplus}\C g.
\]
The Hardy space controls all eigenvectors over the open disk, so the resulting
scalar kernel remains sharp.  The identity $zg=g-1$ makes multiplication by
$z$ bounded and gives an explicit triangular $2\times2$ block matrix.

The added one-dimensional quotient leaves a detectable spectral fingerprint:
$M_z^*$ acquires the boundary eigenvalue $1$, and no other boundary
eigenvalue.  A nontrivial rotation moves this exceptional eigenvalue around
the circle.  Since similarity preserves the point spectrum of the adjoint,
$M_z$ cannot be similar to all its rotations.

\section{Construction and proof}

Let $H^2=H^2(\D)$ with Szeg\H{o} kernel
$k(z,w)=(1-z\overline w)^{-1}$, and put
\[
 g(z):=\frac1{1-z}.
\]
Because $g\notin H^2$, every element of the algebraic sum
\[
 \Hh:=H^2\dotplus\C g
\]
has a unique representation $h+cg$.  Make this an orthogonal Hilbert sum:
\[
 \|h+cg\|_{\Hh}^2:=\|h\|_{H^2}^2+|c|^2.
\]
Point evaluation is bounded on each summand.  Thus $\Hh$ is a scalar
reproducing-kernel Hilbert space with kernel
\begin{equation}\label{eq:kernel}
 K(z,w)=\frac1{1-z\overline w}
       +\frac1{(1-z)(1-\overline w)}.
\end{equation}
The kernel is strictly positive definite: its first summand is the strictly
positive-definite Szeg\H{o} kernel, and the second is positive semidefinite.

Let $T=M_z$ on $\Hh$, and identify $\Hh$ unitarily with
$H^2\oplus\C$ by $h+cg\leftrightarrow(h,c)$.  Since $zg=g-1$,
\begin{equation}\label{eq:T}
 T(h,c)=(Sh-c\mathbf1,c),
 \qquad
 T^*(h,c)=(S^*h,c-\langle h,\mathbf1\rangle),
\end{equation}
where $S=M_z$ on $H^2$ and $\mathbf1$ is the constant function.  In
particular, $T$ is bounded.  Its upper-triangular block form has diagonal
entries $S$ and $1$, hence
\begin{equation}\label{eq:spectrum}
 \sigma(T)=\overline\D.
\end{equation}

\begin{lemma}\label{lem:pointspectrum}
The point spectrum of $T^*$ is
\[
 \sigma_{\mathrm p}(T^*)=\D\cup\{1\}.
\]
Moreover, for every $w\in\D$,
\[
 \ker(T^*-\overline w I)=\spn\{K(\,\cdot\,,w)\}.
\]
\end{lemma}

\begin{proof}
Write $k_w(z)=(1-z\overline w)^{-1}$.  The standard shift calculation gives
$S^*k_w=\overline w k_w$, and each eigenspace of $S^*$ over $\D$ is
one-dimensional.  Solving
\[
 T^*(h,c)=\overline w(h,c)
\]
in \eqref{eq:T} first gives $h=a k_w$.  Since
$\langle k_w,\mathbf1\rangle=1$, the scalar equation then gives
\[
 c=\frac{a}{1-\overline w}.
\]
Under the identification $\Hh=H^2\oplus\C$, the reproducing vector from
\eqref{eq:kernel} is precisely
\[
 K(\,\cdot\,,w)\longleftrightarrow
 \left(k_w,\frac1{1-\overline w}\right).
\]
This proves the kernel identity and sharpness over $\D$.

For a general eigenvalue $\lambda$, the first component in \eqref{eq:T}
requires $S^*h=\lambda h$.  If $|\lambda|<1$, the same calculation supplies
one eigenvector.  If $|\lambda|=1$, then $h=0$, and the second component is
$c=\lambda c$; it has a nonzero solution exactly when $\lambda=1$.  There
are no eigenvalues outside $\overline\D$.  This proves the asserted point
spectrum.
\end{proof}

\begin{theorem}\label{thm:main}
The kernel $K$ in \eqref{eq:kernel} is a scalar-valued sharp positive-definite
kernel on $\D$, its coordinate multiplier is bounded, and that multiplier is
not weakly homogeneous.  Consequently the scalar sharp-kernel question in
\cite[Section 3, p.~8]{Ghara} has an affirmative answer.
\end{theorem}

\begin{proof}
Only weak homogeneity remains to be excluded.  Fix
$\theta\notin2\pi\mathbb Z$ and take the disk rotation
$\rho_\theta(z)=e^{i\theta}z$.  Functional calculus gives
$\rho_\theta(T)=e^{i\theta}T$.  Lemma~\ref{lem:pointspectrum} yields
\[
 \sigma_{\mathrm p}\bigl(\rho_\theta(T)^*\bigr)
 =e^{-i\theta}\sigma_{\mathrm p}(T^*)
 =\D\cup\{e^{-i\theta}\},
\]
whereas $\sigma_{\mathrm p}(T^*)=\D\cup\{1\}$.  The two sets differ.

If $T$ and $\rho_\theta(T)$ were similar, their adjoints would also be
similar, and similar operators have identical point spectra.  This
contradiction shows that $T$ is not similar even to this rotation of itself,
so it is not weakly homogeneous.
\end{proof}

\section{Verification and novelty bounds}

The proof is checkable from the two displayed block equations.  In
particular: (i) $g\notin H^2$ makes the direct sum unambiguous; (ii) $zg=g-1$
proves invariance and boundedness; (iii) the complete eigenvalue equations for
$T^*$ prove both sharpness and the exceptional boundary point; and (iv) a
single rotation is enough to refute weak homogeneity.

A bounded novelty search covered the four run indexes, the local arXiv source
corpus, web/arXiv exact-phrase searches for ``scalar-valued sharp kernel'' and
``not weakly homogeneous'', the source title and arXiv id, and searches for
the specific space $H^2\dotplus\C(1-z)^{-1}$.  It found the source paper and
the author's thesis restating the problem, but no later answer or occurrence
of this construction.  This is evidence only, not a guarantee of novelty.

Two other extraction signals in the source do not affect Theorem~\ref{thm:main}.
The headline Bagchi--Misra question is answered in the source itself.  The
stated Shields power-growth conjecture is known to be false, as discussed by
Bonilla and M\"uller \cite{BM}.

\paragraph{Human-review recommendation.}
High priority: the example is short and settles the literal scalar-kernel
question.  A reviewer should confirm that the source intends exactly its
stated definition of ``sharp'', rather than an unstated Cowen--Douglas
surjectivity hypothesis, and should independently repeat the novelty search.

\enlargethispage{2\baselineskip}
\small
\sloppy
\raggedright
\begin{thebibliography}{9}
\bibitem{Ghara}
S. Ghara,
\emph{The orbit of a bounded operator under the M\"obius group modulo
similarity equivalence},
arXiv:1811.05428; Israel J. Math. 238 (2020), 167--207.

\bibitem{BM}
A. Bonilla and V. M\"uller,
\emph{Kreiss bounded and uniformly Kreiss bounded operators},
arXiv:1912.07931.
\end{thebibliography}
\fussy
\normalsize

\end{document}
